\section{Parameters and the road}\label{sec:parameters}

\subsection{Constants}

The protocol is parameterized by a small, compiled-in set. \code{genesis}
selects the deployment's schedule; every other schedule value is frozen. The
grandfathered sat is a separate compiled exception, not a schedule parameter.
The reference deployment uses the values in Table~\ref{tab:constants}.

\begin{table}[H]
\centering
\small
\caption{Protocol constants for wire version \code{0x08}.}
\label{tab:constants}
\begin{tabularx}{\linewidth}{@{}p{0.27\linewidth}p{0.23\linewidth}X@{}}
\toprule
\textbf{Name} & \textbf{Value} & \textbf{Meaning} \\
\midrule
\code{wire version} & \code{0x08} & The current wire version \\
\code{legacy wire version} & \code{0x07} & Grandfathered sat only \\
\code{mint magic} & \code{"FRG"} & Mint envelope prefix \\
\code{mint action} & \code{0x01} & Mint action byte \\
\code{state magic} & \code{"FRG8CROSS"} & 9-byte crossing body prefix \\
\code{legacy state magic} & \code{"FRG7CROSS"} & Grandfathered sat only \\
\code{state body length} & 55 bytes & Fixed crossing body size \\
\code{grandfathered sat} & \footnotesize\code{1,965,983,893,977,665} & One v7 Hopper; \code{0} = none \\
\code{activation height} & 974{,}034 & Frozen inscription cutover \\
\code{direction encoding} & \code{0,1,2} & Left, straight, right \\
\code{wrap function} & \footnotesize\code{u32LE(H)[0..4] mod 3} & The block's one direction \\
\code{crossing leaves} & 3 + 1 send & One leaf per direction, plus send \\
\code{cross_finality} & 6 blocks & Settlement delay after $T$ \\
\code{card_value} & 25{,}000 sats & Committed card value \\
\code{anchor output value} & 1{,}000 sats & Value of a frog-bearing output \\
\code{dust (P2WPKH)} & 546 sats & Dust threshold for change \\
\code{dust (P2TR)} & 330 sats & Dust threshold for Taproot change \\
\code{NUMS internal key} & \multicolumn{2}{@{}p{0.66\linewidth}@{}}{\footnotesize\code{50929b74c1a04954b78b4b6035e97a5e078a5a0f28ec96d547bfee9ace803ac0}} \\
\bottomrule
\end{tabularx}
\end{table}

\subsection{Generations and the road's difficulty}

The required number of crossings runs in generations. A generation starts at
$b$ crossings, adds one every $\mathit{gi}$ blocks, reaches the ceiling $c$,
holds there for $h_{\rm hold}$ blocks, and then resets to $b$. The values are
\begin{equation}
b=5,\qquad \mathit{gi}=144,\qquad c=25,\qquad h_{\rm hold}=3{,}600,
\end{equation}
so the generation length is
\begin{equation}\label{eq:genlen}
L=(c-b)\mathit{gi}+h_{\rm hold}=(25-5)\cdot144+3{,}600=6{,}480.
\end{equation}
Generation $k$ (1-based) starts at
\begin{equation}
S(k)=G+(k-1)L,
\end{equation}
and owns the half-open interval $[S(k),S(k+1))$. Its number at height $h$ is
\begin{equation}\label{eq:generation}
\gamma(h)=
\begin{cases}
1, & h<G,\\[2pt]
1+\left\lfloor\dfrac{h-G}{L}\right\rfloor, & h\ge G,
\end{cases}
\end{equation}
and the requirement is
\begin{equation}\label{eq:required}
r(h)=
\begin{cases}
b, & h<G,\\[2pt]
\min\!\left(b+\left\lfloor\dfrac{h-S(\gamma(h))}{\mathit{gi}}\right\rfloor,\;c\right), & h\ge G.
\end{cases}
\end{equation}

A generation therefore starts at five required crossings, adds one every 144
blocks, reaches twenty-five at $S(k)+2{,}880$, holds twenty-five for 3{,}600
blocks, and restarts at five. The indexer rejects a published parameter block
whose \code{generation_len} is not
$(c-b)\mathit{gi}+h_{\rm hold}$.

The function is a \emph{sawtooth}, not monotonic. The reset at a generation
boundary is deliberate and load-bearing; it must not be ``smoothed''. The
target used by a crossing is $r(T)$, where $T$ is the crossing's confirmation
height. It is frozen when the crossing confirms and does not follow a later
reset.

A Hopper's \emph{generation} is the generation of its mint confirmation
height, $\mathit{gen}=\gamma(\mathit{mint\_height})$. It is immutable for the
frog's life and is what the art renders as \code{G1}, \code{G2}, and so on. A
frog born in G1 can still hop in G3; only its label is fixed, never its
target.

\begin{figure}[htbp]
\centering
\begin{tikzpicture}
\begin{axis}[
  width=0.97\linewidth, height=6.2cm,
  xmin=0, xmax=12960, ymin=0, ymax=27,
  xlabel={Blocks after genesis}, ylabel={Required crossings},
  xtick={0,2160,4320,6480,8640,10800,12960},
  xticklabels={0,{2{,}160},{4{,}320},{6{,}480},{8{,}640},{10{,}800},{12{,}960}},
  scaled x ticks=false,
  ytick={0,5,10,15,20,25},
  grid=major,
  grid style={draw=deepgreen!12},
  axis lines=left,
  axis line style={draw=deepgreen!55},
  tick style={draw=deepgreen!55},
  label style={font=\small\sffamily, color=ink},
  tick label style={font=\scriptsize\sffamily, color=muted},
  legend style={draw=none, fill=none, font=\scriptsize\sffamily}
]
\addplot[deepgreen, line width=1.2pt] table[x=rel,y=req,col sep=comma]{figures/generation.csv};
\addplot[only marks, mark=*, mark size=2.2pt, color=accent] coordinates {
  (0,5) (2880,25) (6480,5) (9360,25) (12960,5)
};
\node[anchor=south west, font=\scriptsize\sffamily, color=midgreen]
  at (axis cs:2940,25.4) {$+2{,}880$};
\node[anchor=north west, font=\scriptsize\sffamily, color=midgreen]
  at (axis cs:6540,4.6) {reset};
\end{axis}
\end{tikzpicture}
\caption{Two generations of the schedule. Each climbs from five to
twenty-five over 2{,}880 blocks, holds the ceiling for 3{,}600 blocks, then
resets. The sawtooth is intentional.}
\label{fig:generation}
\end{figure}

\subsection{The road's wrap}

There is no per-Hopper death table and no per-Hopper randomness. Each block
has exactly one direction, and every crossing that confirmed in the previous
block is judged against it:
\begin{equation}\label{eq:wrap}
w(B)=\LE_{32}\!\left(H_B^{\rm int}[0..4]\right)\bmod 3,
\qquad
\text{dead}\iff \mathit{dir}=w(T+1).
\end{equation}
Here $H_B^{\rm int}$ is the 32-byte block hash in internal byte order --- the
reverse of the hex an explorer or ord returns --- and $\LE_{32}$ reads its
first four bytes as a little-endian \code{u32}. $\mathit{dir}$ is the
direction the player chose, carried in the crossing body
(Section~\ref{sec:inscriptions}). Direction $0$ is left, $1$ straight, $2$
right.

Normative consequences:

\begin{itemize}[leftmargin=1.4em,itemsep=0.3em]
  \item The wrap is \textbf{global}. Every Hopper whose crossing confirmed in
        block $T$ is judged against the same $w(T+1)$. If the block wraps
        right, every frog in that block that hopped right dies and every
        other frog in that block survives.
  \item The wrap is a function of the block hash alone --- not of the sat, the
        owner, the name, the target or the number of frogs in the block ---
        so an ordinary player cannot choose or predict its own outcome. A
        block producer can influence the candidate hash or withhold a block;
        that is miner influence, not player agency.
  \item Since $2^{32}=3\cdot1{,}431{,}655{,}765+1$, the modulo is biased by
        one part in $2^{32}$ toward left. The bias is accepted. An
        implementation must use this exact integer; it must not substitute a
        wider word, a rejection scheme, or a hash of the hash.
\end{itemize}

\subsection{Parameter commitment}

The parameters are committed to as bytes, so two nodes with different
schedules cannot share a snapshot or a state root. The fields, in order, are
\code{genesis}, \code{base}, \code{growth_interval}, \code{ceiling},
\code{hold}, \code{generation_len}, \code{cross_finality} and
\code{card_value}. The canonical encoding is the 36-byte concatenation
\begin{equation}\label{eq:canon}
\begin{aligned}
\operatorname{canon}(P)={}&
\LE_{32}(G)\,\|\,\LE_{32}(b)\,\|\,\LE_{32}(\mathit{gi})\,\|\,
\LE_{32}(c)\\
&{}\|\,\LE_{32}(h_{\rm hold})\,\|\,\LE_{32}(L)\,\|\,
\LE_{32}(6)\,\|\,\LE_{64}(25{,}000),
\end{aligned}
\end{equation}
and the parameter hash is
\begin{equation}\label{eq:paramhash}
H_{\rm params}=\SHA\!\left(\operatorname{canon}(P)\right).
\end{equation}
For the reference deployment,
\begin{center}
\code{d466f8db3efa330dc92444274675804e47ec76b62393eeae88063ef764bcf72f}.
\end{center}
A wallet refuses to mint or hop when the indexer's published $H_{\rm params}$
differs from its own. It also refuses an indexer that does not answer the
parameter gate, or whose \code{network} does not match the wallet's network.
Missing identity is not permission to sign. ord remains the canonical source
for sat ranges and inscription validity; the HOPPER indexer is for discovery
and the leaderboard, never for overriding ord. A snapshot whose stored hash
differs is a hard error, never a migration.

\subsection{Outcome timing}

$w(T+1)$ is learned when block $T+1$ is mined; the outcome is applied at
$T+\kappa=T+6$. Nothing about a Hopper's state changes in between, and a
reorg that replaces $T+1$ replaces the wrap. The wrap is recorded once per
block and applied to every crossing in the previous block, so the game is
decided by the chain at block granularity, not by a per-frog draw.
